\section{Measurement, Decoherence, and the Emergence of Classicality}
  \label{sec:classical-limit}

  The structural origin of Bell-type correlations identified above naturally raises
  the question of their apparent absence in classical macroscopic phenomena.
  Everyday physical systems exhibit effectively local and statistically independent
  behavior, despite being ultimately composed of quantum constituents.
  Any satisfactory account of Bell violations must therefore explain not only their
  existence, but also their suppression in the classical limit.

  Within the present framework, this transition is understood as a change in the
  effective properties of the projection from underlying configurations to observable
  descriptions.
  While the projection is generally non-injective at microscopic scales, the effective
  descriptive capacity associated with observable states becomes large compared to the
  set of configurations that are actually explored.
  In this regime, the mapping becomes approximately injective.

  Decoherence plays an important but secondary role in this process.
  Through continuous interaction with uncontrolled environmental degrees of freedom,
  the system is dynamically driven toward states that are stable under coarse-grained
  observation.
  This evolution effectively reduces the overlap between distinct underlying
  configurations contributing to the same observable outcome, thereby pushing the
  system toward a regime in which the projection behaves as approximately one-to-one.
  In this regime, the multiplicity of pre-images associated with each observable outcome
  becomes negligible, and the projection entropy effectively vanishes.

  From the perspective of observable descriptions, this corresponds to a regime in
  which the effective capacity of the description exceeds the occupied portion of
  configuration space.
  Global consistency constraints associated with non-injective projections become
  negligible, and observable outcomes can be treated as effectively independent.

  In this regime, joint probabilities admit an approximate decomposition of the form
  \begin{equation}
    P(a,b \mid A,B) \approx P(a \mid A)\,P(b \mid B),
  \end{equation}
  up to corrections that are exponentially small in the strength of environmental
  coupling or the degree of coarse-graining.
  Bell inequalities are therefore not violated in practice, not because the underlying
  structure has changed, but because the effective description operates in an
  approximately injective regime.

  Importantly, this recovery of classical behavior does not require the introduction
  of additional collapse postulates or modifications of quantum dynamics.
  It follows from the same structural mechanism responsible for Bell violations,
  applied in a regime where the effective descriptive capacity is large compared to
  the occupied state space.
  Classicality thus corresponds to a limit in which observable descriptions admit an
  approximately complete specification, rendering joint probabilities effectively
  factorizable.

  From this viewpoint, the quantum-to-classical transition reflects a change in
  descriptive accessibility rather than a change in fundamental physical laws.
  Bell-type correlations are suppressed when finite-capacity effects become negligible
  in the relevant regime.
  The classical world emerges as a limit in which effective injectivity is restored,
  and with it, the validity of classical notions of separability and locality.
